% ============================================================================
% DRAFT reframed Introduction v2 (Vikas theme 2026-07-17 + Codex review + verified
% cites). Wide-band EMPIRICAL SHAPE HUNT of single pulses, TEMPORAL + SPECTRAL.
% Replaces sec:intro AFTER approval. NEW bib keys needed are listed at the bottom
% (all verified 2026-07-17). [CITE?] = still to add/confirm.
% ============================================================================
\section{Introduction} \label{sec:intro}

Gamma-ray bursts (GRBs) release most of their energy as a brief, highly variable
flash of $\sim$keV--MeV photons---the prompt emission---whose spectrum evolves
strongly on sub-second timescales \citep{Norris1996,Kaneko2006,Lu2012,Yu2016}.
To first order this spectrum is described by empirical functions such as the Band
function \citep{Band1993} or a cutoff power law: smooth, curved forms peaking in
$\nu F_\nu$ (the energy-times-flux representation, whose maximum \Ep\ marks where
the burst radiates most of its power), with a low-energy photon index $\alpha$.
These functions fit well but are physically agnostic---the same curvature can
arise from very different mechanisms \citep{KumarZhang2015}---and after three
decades the radiative origin of the prompt emission is still debated.

Much of the apparent complexity is temporal blending. \citet{HakkilaPreece2011}
showed that the ensemble ``intensity-tracking'' behavior of complex bursts, in
which \Ep\ appears to follow the flux, is naturally reproduced by the
\emph{superposition} of two or more simpler hard-to-soft pulses; isolated pulses,
by contrast, are dominated by clean hard-to-soft evolution. To recover the
\emph{intrinsic} spectral behavior without this contamination, we restrict our
analysis to single-pulse bursts---one non-overlapping rise--peak--decay, long the
cleanest laboratory for prompt-emission studies
\citep{RydePeer2009,Basak2014,Yu2016,Busby2024}.

No single spectral shape has proven adequate even for these clean pulses, and the
literature offers a zoo of alternatives. Optically-thin synchrotron from a cooling
electron population radiates between two characteristic frequencies and so breaks
at \emph{two} energies \citep{Sari1998}, motivating double smoothly-broken power
laws that recover the expected $-2/3$ and $-3/2$ slopes in bright bursts
\citep{Oganesyan2017,Oganesyan2018,Ravasio2018,Ravasio2019}; indeed, a physical
synchrotron model with time-dependent electron cooling fits $\sim$95\% of the
time-resolved single-pulse GBM spectra it was applied to
\citep{Burgess2019cool}.\footnote{J.~M. Burgess et al., \textit{Nat. Astron.}
2019 (arXiv:1810.06965).} Other pulses require a quasi-thermal component from the
jet photosphere---a blackbody superposed on a non-thermal continuum
\citep{Ryde2004,Guiriec2010,Guiriec2011,Burgess2014a}---occasionally as a distinct
\emph{third} spectral component \citep{Guiriec2015}, and still other continua or
cutoffs have been invoked, produced for instance by a fast-cooling spectrum with an
exponential cutoff above the peak \citep{Varun2025}. That such competing shapes can
coexist in a single, clean pulse motivates fitting them all and letting the data
decide, rather than adopting any one a priori.

Discriminating among these shapes demands the widest possible spectral coverage. A
low-energy break is easily masked from below, and a cutoff or an additional
high-energy component sits above the peak, near or beyond the reach of the Gamma-ray
Burst Monitor (GBM; $8\,\mathrm{keV}$--$40\,\mathrm{MeV}$; \citealt{Meegan2009}).
Fermi uniquely extends the band by a decade through the LAT Low-Energy events
(LLE; $30$--$100\,\mathrm{MeV}$) and the LAT itself
($>100\,\mathrm{MeV}$; \citealt{Atwood2009}), and high-energy cutoffs are directly
seen in bursts bright enough to register there. Wherever a pulse is bright enough,
we fold these data into a joint fit, giving coverage from $8\,\mathrm{keV}$ to above
$100\,\mathrm{MeV}$---the widest band the data allow.

A single pulse is equally clean in the \emph{time} domain, and we treat its temporal
and spectral properties as two views of the same event rather than separate studies.
Alongside the spectra we characterize the pulse duration and hardness, the minimum
variability timescale and its dependence on energy \citep{Golkhou2014,Golkhou2015},
the spectral lag between energy bands, and the pulse shape itself---its rise
timescale, curvature at the peak, and decay slope
\citep{Norris2005,Kocevski2003,Gowri2025}---and we ask directly how the spectrum
during the rise, at the peak, and during the decay relates to that morphology.

Our approach throughout is deliberately empirical and model-agnostic. Rather than
assume a radiative process, we fit a broad, pre-specified menu of empirical
functions---power laws with one or two breaks, with free or fixed peak curvature,
with and without an added blackbody or high-energy component---uniformly to
\emph{every} time-resolved spectrum, and let the data select the preferred shape. We
reduce every burst identically with \textsc{3ML} \citep{Vianello2015} on
Bayesian-block time bins, which follow the spectral evolution more faithfully than
fixed signal-to-noise bins \citep{Scargle2013}. This work differs from previous
prompt-spectral catalogs in combining a reproducible single-pulse morphology
selection \citep{Busby2024} with a pre-specified hierarchy of empirical shape
families on common time-resolved intervals, and in adding a separately-defined joint
GBM--LLE/LAT analysis for high-energy structure.

We therefore aim at an empirical census of shapes and their evolution, asking, at
scale: (i) which spectral shapes single pulses take, and how often; (ii) how the
parameters and the preferred shape evolve through a pulse; and (iii) which
correlations emerge---among them the hard-to-soft trend \citep{Lu2012,Busby2024}
and, where the data require it, a relation between two spectral breaks or an
\Ep--\kT\ relation measured from a single composite fit. We stress that two-break,
thermal, and cutoff signatures are outcomes to be \emph{measured}, not assumptions.
The paper is organized as follows. Section~\ref{sec:method} describes the sample,
data reduction, and spectral and temporal analyses; Section~\ref{sec:obs} presents
the shape census, parameter distributions, evolution, and correlations;
Section~\ref{sec:discussion} discusses the results; and Section~\ref{sec:summary}
summarizes. Uncertainties are quoted at the $1\sigma$ (68\%) level unless stated
otherwise.

% ============================================================================
% NEW bib entries to add to two_break.bib (all verified 2026-07-17):
%  Burgess2019cool : Burgess, Begue, Bacelj, Giannios, Berlato & Greiner 2019,
%                    Nat. Astron., "Gamma-ray bursts as cool synchrotron sources",
%                    arXiv:1810.06965  (single-pulse GBM; ~95% of the spectra)
%  HakkilaPreece2011: Hakkila & Preece 2011, ApJ 740, 104 (superposition -> tracking) [CONFIRM vol/page]
%  Varun2025       : Varun et al. 2025, GRB 241030A, arXiv:2510.24864 (SBPLxCut, -2/3,-3/2 + cutoff)
%  Norris2005      : Norris et al. 2005, ApJ 627, 324 (FRED pulse function)
%  Kocevski2003    : Kocevski, Ryde & Liang 2003, ApJ 596, 389 (KRL asymmetric pulse)
%  Gowri2025       : Gowri, Pe'er, Ryde & Dereli-Begue 2025, ApJ 991, 230, arXiv:2409.17860 (two-sigmoid pulse)
%  Golkhou2014     : Golkhou & Butler 2014, ApJ 787, 90 (MVT via Haar wavelets)
%  Golkhou2015     : Golkhou, Butler & Littlejohns 2015, ApJ 811, 93 (MVT vs energy, 938 GBM)
% DISCUSSION-only (do not put in intro): Ghisellini et al. 2020, arXiv:1912.02185 (proton synchrotron)
% REMOVED misuses: Ravasio2023 (=221009A MeV emission line, NOT narrow peaks); the vague
%   Amati "second-component energetics" aim; Preece1998 for superposition; Ronchi2020 for proton-synch.
% ============================================================================
